% Exact low-depth historical parameters, separately from the general E.33.
\subsection{Exact historical class-three constructions and quantifiers}

The general fixed-depth theorem E.33 does not by itself verify
the smaller constants in the earlier class-three constructions.
We now keep those constants. All finite reductions below are
reductions of integral class-two sections unless a different
domain is explicitly specified.

\begin{result}{E.60}{Proved}{The scaled-factorial first-prime exclusion}
\label{coh:scaled-five}
There is an ordinary infinite-support pair
\[
 A=3\prod_{j\geq0}q_j^{5(j+4)!},\qquad
 B=-2\prod_{j\geq0}r_j^{5(j+4)!}
\]
whose integral class-two fibre consists of exactly twenty-four
geometric-conjugacy classes, every one of which has nonzero
third obstruction at $q_0$ in coefficient five, already
modulo $25$. Thus its first empty unrestricted integral
nilpotent fibre is class three. This confirms the assertion
and exact exponent sequence of GT C.2.5.
\end{result}
\begin{proof}
Use the actual point-three field $L_5$ of (E.37.0).
It has group $H_5\rtimes\mathbb F_5^\times$ and maximal
abelian subfield $\mathbb Q(\mu_5)$, by the proof there.
Put
\[
 n_j=5(j+4)!,\quad
 B(n)=(n+1)^4,\quad
 H_j(p)=2v_p(n_{j+1})+j+3 .
\]
At the first stage $n_0=120,n_1=600$. Choose the exact
old local points $(1,1)$ of $3X^{120}-2Y^{120}=1$.
Thus the first $q$-CRT class is one modulo
\[
 M=\prod_{p\leq B(600)}p^{H_1(p)} .
\]
The primes dividing the initial coefficients are already
in this product. Set
\[
 D=\operatorname{lcm}(M,120,25),\quad
 F=\mathbb Q(\mu_D),\quad
 R=F((-2)^{1/120},3^{1/5}).
\]
The extension $R/F$ is abelian. On the other hand
$L_5\cap F=\mathbb Q(\mu_5)$, so
$\operatorname{Gal}(L_5F/F)=H_5$.
Its intersection with $R$ is abelian over $F$ and
is therefore fixed by the center
$[H_5,H_5]$. Any chosen nonzero central element,
paired with the identity on $R$, consequently defines
an automorphism of the actual compositum.
Chebotarev \cite[8.31, p.147]{CohMilneANT} gives
arbitrarily large primes $q_0$ satisfying all the
first-stage conditions and
\[
 q_0\equiv1\pmod{25},\qquad
 \rho_{3,5}(\operatorname{Frob}_{q_0})
                  =(0,0,c;1),\quad c\ne0.
 \tag{E.60.1}
\]
The property of being nonzero central is invariant
under conjugation in the whole semidirect product.
Choose the remaining primes by the original two-prime
recursion of E.102 with the displayed scaled exponents.

Here are all parameter changes needed to apply that
arithmetic proof. Every $n_j$ is divisible by $24$,
$n_j\mid n_{j+1}$, the ratio is still $j+5$, and
every prime dividing $n_{j+1}$ is at most $B(n_j)$.
The coefficient forms $3u^{24},-2v^{24}$ in particular
retain the forms $3u'^8,-2v'^8$ used for the radical
intersection. The old-prime precision
$H_j(p)>2v_p(n_{j+1})$ is unchanged in form, so the
same sharpened Hensel estimate proves the next unit
Fermat conditions. The finite-field bound is still
$B(n)$, applied to the actual new exponent.
Finally $n_j$ is eventually divisible by every
integer. Hence the products converge in every global
power quotient, their valuations stabilize at every
finite prime, and their local units converge.
The increasing inequalities
$v_p(a_j+b_j-1)\geq H_j(p)$ make the local sums
equal to one. Their nonzero valuation patterns are
$(1,0)$ at three, $(0,1)$ at two, $(n_j,0)$ at
$q_j$, and $(0,n_j)$ at $r_j$. This proves
ordinarity and infinite support with these exact
exponents.

Align any integral class-two section with the literal
Kummer sums of these products. Its abelian coordinates
agree with those of point three modulo five.
The difference of the central coordinates is therefore
a $\mathbb F_5(2)$-cocycle. At a fresh prime, the
actual integral tame relation gives central inertia
$n_ju$ or $-n_jv$, with $u,v\in\mathbb Z_5$;
it is zero modulo five. At zero-valuation primes
the same relation gives central inertia zero
integrally. At two and three the difference on
inertia is killed by $p(p-1)$, a five-adic unit.
Thus the difference is unramified outside five
and is a coboundary by (E.37.1).
At $q=q_0$, the cyclotomic character is one modulo
five, so that coboundary vanishes on Frobenius.
In the notation of B.23, every integral lift has
\[
 \alpha\equiv0,\qquad\gamma\equiv c\ne0\pmod5.
\]
Here $m=n_0=120$, $n=0$, $\beta=(q-1)u$, and
$t=mu$. The exact positive-support substitution in
B.23 is
\[
 Z_{D,q}=m(\gamma-q\alpha u)
          +(q^2-1)u\binom m2 .
 \tag{E.60.2}
\]
Its first term has five-adic valuation one, and
its second term has valuation at least two.
Thus $v_5(Z_{D,q})=1<v_5(q^2-1)$.
B.22 identifies this with the actual nonzero local
third obstruction, whose reduction modulo $25$
is nonzero as well. A.25 and A.27 give existence
and the twenty-four classes at depth two.

The prime rule is recursive without using a generic
radical equation for the Heisenberg cover. One may
fix $L_5$ by enumerating marked finite Galois
extensions with group $H_5\rtimes\mathbb F_5^\times$,
the specified cyclotomic and Kummer characters,
and ramification only at $2,3,5$.
Such a marked extension exists by point three.
Any two have central difference unramified outside
five: the only possible extra places are two
and three, where $p(p-1)$ is a unit. Equation
(E.37.1) makes them geometrically conjugate.
Finite number-field operations verify the marking,
ramification, and Frobenius conditions. Having
fixed one such field, take the least eligible prime
at each step and the least eligible finite residue
choices. This confirms the asserted recursive
construction. It does not certify the optional
generic $\Theta$ equation in the historical proof,
which is not required here.
\end{proof}

Let $\mathscr A$ be the space of original E.102
prime sequences, with $n_j=(j+4)!$, including
their legal finite CRT choices, and give it the
prefix topology. Let $\mathscr A_c$ be the
subspace passing all the compact elliptic support
tests in E.113. Let $\mathscr F_3$ consist of
outputs whose first empty integral fibre is class
three.

\begin{result}{E.61}{Proved}{Class-three failure after every original finite prefix}
\label{coh:exact-dense-three}
For every $k\geq0$,
\[
 \operatorname{pr}_k(\mathscr F_3)
     =\operatorname{pr}_k(\mathscr A),\qquad
 \operatorname{pr}_k(\mathscr F_3\cap\mathscr A_c)
     =\operatorname{pr}_k(\mathscr A_c).
\]
One may take a coefficient prime $\ell>k+4$,
$\ell\geq5$, and the distinguished later support
prime $q_{\ell-4}$. Every integral class-two
section then has nonzero local third obstruction
there, detected modulo $\ell^2$. All exponents
remain exactly $(j+4)!$. This confirms GT C.2.8.
\end{result}
\begin{proof}
Choose $\ell$ outside the finite old support and put
$J=\ell-4$. Fill all intervening stages before
$J$ by the original rules, avoiding $\ell$.
For the compact-passing assertion, impose the
additional noncube choices of E.114 at those
new stages. These are compatible with every
old legal CRT class, as proved there.
Write $a=a_J,b=b_J$, and let $S_0$ consist of
the fresh support primes before $J$.
Then
\[
 n=n_J=\ell!,\quad m=n_{J+1}=(\ell+1)!,
 \quad v_\ell(n)=1,\quad \ell\nmid ab,\quad
 \ell\nmid n_i\quad(i<J).
 \tag{E.61.1}
\]
At the stage-$J$ local coefficient prime take
$\xi_\ell=1$. The original invariant and the
sharpened Hensel estimate give
\[
 v_\ell(a+b-1)\geq H_J(\ell)=\ell+1,\quad
 v_\ell(\eta_\ell-1)\geq\ell,\quad
 v_\ell(M)=H_{J+1}(\ell)=\ell+2.
 \tag{E.61.2}
\]
Thus the $q_J$-CRT class $Q_0$ is one modulo
$\ell^2$, as well as one modulo $m$.

We must freeze the actual finite class-two
representation before choosing $q_J$. Choose
positive rational integer placeholders $P,R$,
prime to $S_0$, with
$P\equiv\xi_p,\ R\equiv\eta_p\pmod p$
at every $p\in S_0$. Fix integral Kummer cocycles
and put
\[
 a^*=\kappa_\ell(a)+n\kappa_\ell(P),\qquad
 b^*=\kappa_\ell(b)+n\kappa_\ell(R).
\]
At an old $A$-support prime $bR^n$ is a
principal unit, since the chosen unit Fermat
point gives $b\eta_p^n\equiv1\pmod p$.
Its integral Kummer cocycle therefore has
a unique integral local primitive $u_p^*$.
At an old $B$-support prime the corresponding
primitive $v_p^*$ is supplied by $aP^n$.
The uniqueness follows from
$H^0(G_p,\mathbb Z_\ell(1))=0$.
Prescribe central inertia modulo $\ell$ to be
$n_i u_p^*$ or $-n_i v_p^*$, respectively.

The finite cup of $(a,b)$ modulo $\ell$ is
zero. At every finite place the invariant gives
a point of $aX^n+bY^n=1$ with $n$ divisible
by $\ell$, so its Kummer cup vanishes locally.
Finite weight-two detection in A.25 makes it
zero globally. Start with any finite primitive.
The prescribed old inertia values are compatible
with the same finite abelian cocycles, since
$(a^*,b^*)\equiv(\kappa(a),\kappa(b))\pmod\ell$.
Prescribe zero inertia elsewhere outside
$S_0\cup\{2,3,\ell\}$.
The differences from the initial primitive are
a finitely supported vector of residue classes.
The finite residue surjection in A.25 supplies
a global cocycle correcting them.
This constructs a primitive $h_*$ with those
inertia restrictions. Its ambiguity is a
central coboundary, by (E.37.1); at two and
three no extra inertia ambiguity remains because
$p(p-1)$ is a unit modulo $\ell$.

Require every later prime, after stage $J$, to
satisfy
\[
 q_i\equiv r_i\equiv1\pmod{\prod_{p\in S_0}p}
                  \qquad(i>J).
 \tag{E.61.3}
\]
For any such continuation, choose the literal
abelian cocycles as the convergent sums of
its rational Kummer factors. At an old $p$,
$q_J/P$, $r_J/R$, and all later factors are
principal units. Divide the cocycle differences
from $a^*,b^*$ by $\ell$. Each summand is an
integral local coboundary. Its coefficient
$n_i/\ell$ tends to zero, so the corresponding
zero-cochain series converges. Hence
\[
 a_\infty-a^*=\ell\delta e_p,\qquad
 b_\infty-b^*=\ell\delta f_p
                       \quad(p\in S_0).
 \tag{E.61.4}
\]
The old local primitives, and therefore the old
central inertia coordinates, agree modulo $\ell$.
At all new support primes the exponents are
divisible by $\ell$, so integral central inertia
is zero modulo $\ell$; at zero-support primes
it is zero integrally.
Any reduced integral lift thus differs from
$h_*$ by a cocycle unramified outside $\ell$,
which is a coboundary by (E.37.1).
This comparison uses division only before
finite reduction.

The finite primitive defines
\[
 \rho_*:G_{\mathbb Q}\longrightarrow
 H_\ell\rtimes\mathbb F_\ell^\times .
 \qquad L_*=\overline{\mathbb Q}^{\ker\rho_*}.
\]
Its image is the whole group. The valuations
one at three and two prove independence of
the two Kummer characters over $\mathbb Q(\mu_\ell)$;
commutators generate the center of $H_\ell$.
As in (E.37.0), the full derived subgroup
is $H_\ell$ and
$L_*^{\rm ab}=\mathbb Q(\mu_\ell)$.

Now put
\[
 F=\mathbb Q(\mu_M),\qquad
 R_b=F(b^{1/n}),\qquad
 R'=R_b(a^{1/\ell}).
\]
The original radical intersection E.102 lets
the prescribed CRT automorphism $\sigma_F$
extend to $R_b$, fixing $b^{1/n}$.
The rational $\ell$-power classes of $a,b$
remain independent in $F$. Indeed any nontrivial
combination has a valuation at two or three
not divisible by $\ell$. Its real $\ell$th
root generates a nonnormal degree-$\ell$
field, which cannot be contained in an
abelian extension of $\mathbb Q$.
The cyclic extension $R_b/F$ has its unique
degree-$\ell$ subfield $F(b^{1/\ell})$,
so it is disjoint from $F(a^{1/\ell})$.
Consequently $\sigma_F$ extends to $R'$
fixing both chosen roots. Moreover
\[
 L_*F\cap R'=F(a^{1/\ell},b^{1/\ell}),
 \qquad [L_*R':R']=\ell .
 \tag{E.61.5}
\]
The intersection is contained in the maximal
abelian extension of $F$ within $L_*F$,
because $R'/F$ is abelian; it contains that
extension because both Kummer roots were adjoined.
This proves equality, not just an inclusion.
A nonzero central element of $H_\ell$ on
$L_*$ and $\sigma_F$ on $F$ agree with the
chosen action on $R'$ along (E.61.5).
They define an automorphism of the actual
compositum.

For the compact-passing version also adjoin
$F(3^{1/3})$. It is disjoint from $R_b$:
the unique cubic subfield of $R_b$ is
$F(b^{1/3})=F((-2)^{1/3})$, and the two
rational cubic classes are independent in a
cyclotomic field by the preceding nonnormal
radical argument. The two subsequent
degree-$\ell$ extensions cannot introduce a
cubic intersection. Thus the chosen compositum
automorphism may also act nontrivially on
$3^{1/3}$. Chebotarev now supplies arbitrarily
large legal $q_J$ with
\[
 q_J\equiv1\pmod{\ell^2},\qquad
 \rho_*(\operatorname{Frob}_{q_J})
                     =(0,0,c;1),\quad c\ne0,
 \tag{E.61.6}
\]
and with three a noncube when required.
Choose $r_J$ by the original rules, retaining
the independent noncube-of-two condition in
that case.

Condition (E.61.3) remains available. The
stage-$J$ CRT equations made the coefficient
sum one modulo every prime of $S_0$.
At a later stage, if such a prime does not
divide the current exponent, the unit
coefficient supplies a simple root near one.
If it divides the exponent, it lies within
the range of the strong invariant, and the
precision $H_i(p)>2v_p(n_{i+1})$ supplies
the needed near-one root, including the next
exponent's congruence. Taking these roots
in CRT proves (E.61.3) while preserving
every other rule. Thus a legal infinite
continuation exists with the original prefix.

At the selected $q$, every integral lift has
$\alpha\equiv0,\gamma\equiv c\ne0\pmod\ell$
by (E.61.4) and (E.37.1).
In (E.60.2) use $m_0=\ell!$.
The first term has valuation one, while
the second has valuation at least two.
The local third obstruction is therefore
nonzero already modulo $\ell^2$.
Class-two existence gives first failure
exactly three. This proves both prefix
equalities. In the compact-passing case,
E.117 also supplies the affine abelian
section with proper class zero; the actual
quotient comparison in E.118 excludes its
metabelian lifts.
\end{proof}

\begin{result}{E.62}{Proved}{Open dense failure and finite forcing prefixes}
\label{coh:exact-open-three}
Let $\mathscr U_3$ be the original E.102
outputs whose local third obstruction is
nonzero at some coefficient $\ell\geq5$
and is detected at a finite exponent.
Then $\mathscr U_3$ is open and dense in
$\mathscr A$, and
$\mathscr U_3\cap\mathscr A_c$ is open
and dense in $\mathscr A_c$.
Every finite legal prefix has a finite
extension all of whose legal ordinary
continuations have first integral failure
at class three. The same statement holds
relative to the compact-passing domain.
This confirms GT C.1.21 and C.2.9.
\end{result}
\begin{proof}
For completeness the exact original
prefix indices, already covered by B.20,
are as follows. Given $\ell\geq5,N\geq1$,
let $J_N$ be least with
$v_\ell((J_N+4)!)\geq N$, and set
\[
 S=\bigl(\{2,3\}\cup
          \{q_i,r_i:i<J_N\}\bigr)\setminus\{\ell\},
 \qquad D=N+\max_{p\in S}v_\ell(p-1).
 \tag{E.62.1}
\]
Extend through $L-1$ with $L\geq J_N$ and
$v_\ell((L+4)!)\geq D$. Any two continuations
have abelian cocycle differences divisible
by $\ell^D$. At a prime in $S$, the
local primitive difference is divisible by
$\ell^{D-v_\ell(p-1)}$, hence by $\ell^N$.
Their integral central inertia coordinates
therefore agree modulo $\ell^N$.
Outside $S$, all nonzero support valuations
are divisible by $\ell^N$, and central
inertia vanishes modulo that power.
The reduced central difference is
unramified outside $\ell$.
For this entire finite exponent the
coefficient sequence and the adjacent
integral groups in (E.37.1) give
$H^1(G_{\mathbb Q,\{\ell,\infty\}},
\mathbb Z/\ell^N(2))=0$.
Thus the reductions are geometrically
conjugate, centrally after aligning the
literal abelian cocycles.

The actual integral class-three Hall laws,
including the arithmetic action, descend
modulo $\ell^N$ for $\ell\geq5$.
Their fixed denominators divide powers
of $24$, which are units here.
The zero-last-coordinate set lift commutes
with reduction, so its finite defect is
the reduction of the integral defect.
A class-two geometric conjugator lifts
through the class-three group. Conjugation
acts trivially on the central weight-three
kernel, so it changes the resulting
defect only by a coboundary.
Consequently the finite obstruction class
is unchanged by the conjugacies just proved.

For an output in $\mathscr U_3$, choose
a detecting finite exponent and use
(E.62.1), also including the detecting
support prime in the prefix. Every
continuation then has the same nonzero
finite local class. This proves openness.
E.61 supplies an output in $\mathscr U_3$
after every legal prefix, proving density.
Applying openness to that output gives
the finite forcing extension. The same
argument with the relative prefix space
gives the compact-passing statements.

At $\ell\geq5$ the global integral class-two
class is unique up to geometric conjugacy
by A.26, so the obstruction excludes all
integral class-two lifts. A.25 makes the
first failure exactly three.
The map from prime sequences to completed
pairs is continuous, since a fixed power
quotient depends on only finitely many
factorial exponents. Hence the density
assertions also hold in the corresponding
images in the power-completion topology.
\end{proof}

The complement of $\mathscr U_3$ is closed
with empty interior, including relatively
inside $\mathscr A_c$. These statements do
not determine whether that complement is
empty. Vanishing of its finite local tests
still leaves the strict quotient-convergence
condition and the higher tower conditions.

\begin{result}{E.63}{Proved}{The exact factorial finite-third-layer theorem}
\label{coh:finite-third-factorial}
Run the original two-prime recursion with
$n_j=(K+j)!$, $K\geq4$, and put $N=K!$.
Let $s_*$ be the unique integral class-two
class selected by the two exceptional-prime
tests of B.101. Its actual third obstruction
is zero in
\[
 H^2(G_{\mathbb Q},
           (\widehat{\mathbb Z}/N\widehat{\mathbb Z})(3))^2.
\]
Equivalently, it lifts to
\[
 E_{3,N}=
 E_3/N(\Gamma_3\Pi/\Gamma_4\Pi)
                 \longrightarrow E_2 .
\]
The full class-two quotient is retained
integrally. This confirms SECTIONS A.52,
including the original $K=4,N=24$ case.
For every fixed positive $M$, some ordinary
infinite-support pair has a lift to $E_{3,M}$.
\end{result}
\begin{proof}
The exact initial-factorial change satisfies
every arithmetic condition in E.102:
$24\mid n_j$, $n_j\mid n_{j+1}$, every
prime of $n_{j+1}$ is at most $K+j+1<B(n_j)$,
and the initial sum $3-2=1$ satisfies
every precision. The same field, Hensel,
finite-field, and limit arguments apply.
In particular, each support prime is one
modulo $n_{j+1}$ and hence modulo $24$.
Define completed positive classes
\[
 A_1=\prod_jq_j^{n_j/N},\qquad
 B_1=\prod_jr_j^{n_j/N}.
\]
Their products converge because $n_j/N$
is eventually divisible by every integer.
In additive Kummer notation,
$A=\kappa(3)+NA_1$, $B=\kappa(-2)+NB_1$.
Modulo two, each root class is represented
by a positive finite product of its support
primes, all congruent to one modulo $24$.

Choose point-three cocycles $a_0,b_0,h_0$.
At a fixed complex conjugation $\sigma$,
choose the real roots so that
$a_0(\sigma)=0$, $b_0(\sigma)=-1$.
The positive classes $A_1,B_1$ admit
literal cocycles $a_1,b_1$ with
$a_1(\sigma)=b_1(\sigma)=0$.
Use
\[
 a=a_0+Na_1,\qquad b=b_0+Nb_1
\]
and any integral central coordinate $h$
supplied by A.25. The cocycle identity
at $\sigma^2=1$ gives $2h(\sigma)=0$,
because $a(\sigma)=0$; hence
$h(\sigma)=h_0(\sigma)=0$ in the
torsion-free integral coefficient module.

For $\ell\mid N$, put $t=v_\ell(N)$.
The difference $c=(h-h_0)\bmod\ell^t$
is a finite weight-two cocycle.
At a fresh support prime its inertia
is $n_ju$ or $-n_jv$, and thus divisible
by $N$. At every other prime outside
$2,3,\ell$ both integral inertia
coordinates are zero.
At two, for $\ell\ne2$, the classes of
$3$ and $A_1$ are locally zero in
$\ell$-adic Kummer cohomology: every
two-adic unit is uniquely $\ell$-divisible
up to its two-primary finite factor,
which also has zero $\ell$-completion.
Write $a_0=\delta u_0,a_1=\delta u_1$.
The central inertia difference is
$-Nu_1$, since the $B$-valuation is one.
At three, for $\ell\ne3$, both $-2$ and
$B_1$ belong to the completed principal
unit subgroup: every $r_j\equiv1\pmod3$.
Their $\ell$-adic Kummer classes vanish.
Writing $b_0=\delta v_0,b_1=\delta v_1$,
the central inertia difference is $Nv_1$.
This direct divisibility avoids losing
a factor of two at the prime three.
Thus $c$ is unramified outside $\ell$
and is zero on complex conjugation.

For odd $\ell$, the fixed-ramification
group $H^2(G_{\mathbb Q,\{\ell,\infty\}},
\mathbb Z_\ell(2))$ is zero by A.25.
The finite class of $c$ therefore lifts
to integral first cohomology.
At $\ell=2$, its Bockstein has zero
real restriction; A.25 proves that
integral fixed-ramification second
cohomology injects into the real group.
Its Bockstein is again zero, so $c$
lifts integrally there too.
A.26 identifies those integral classes
as multiples of $\eta$ at two and three,
and as zero elsewhere. A single
$k\in\mathbb Z/24$ realizes the necessary
two-primary and three-primary twists.
After that twist and an integral central
coboundary we can arrange
\[
                        h\equiv h_0\pmod N .
 \tag{E.63.1}
\]
The literal integral abelian coordinates
are unchanged by this central adjustment.

For odd $\ell\mid N$, the actual two
class-three cocycles in B.23 now agree
with those of point three modulo
$\ell^t$, since their only binomial
denominator two is a unit.
At two, $t=v_2(N)\geq3$, and the exact
binomial calculation is
\[
 \binom{a_0+Na_1}{2}-\binom{a_0}{2}
       \equiv-\frac N2a_1\pmod{2^t}.
 \tag{E.63.2}
\]
Indeed the omitted terms are
$Na_0a_1+N^2a_1^2/2$, both divisible
by $2^t$. The same formula holds
for $b$. Every other term of the
actual defects agrees by (E.63.1).
Their differences are therefore the
images of the finite classes
\[
 (A_1\bmod2)\cup\kappa_2(-2),\qquad
 -(B_1\bmod2)\cup\kappa_2(3)
 \tag{E.63.3}
\]
under the coefficient map
$\mathbb F_2\to\mathbb Z/2^t(3)$
sending one to $N/2$. This map is
equivariant, because its image has
order two.

Both finite classes in (E.63.3) vanish.
Use the finite positive rational
square-class representatives described
above. At two they are one modulo
eight and are squares. At three they
are one modulo three and are squares.
At each of their support primes,
congruence one modulo $24$ makes
both $-2$ and $3$ local squares,
since that prime splits in
$\mathbb Q(\mu_{24})$, which contains
their square roots. At every remaining
odd prime both factors are units,
so the cup is inflated from the
finite residue field and vanishes.
At the real place the first factor
is positive. Finite Brauer localization,
or the finite weight-two injection
proved in A.25, gives global vanishing.
This argument is at finite coefficients;
it asserts no inverse-limit vanishing.
Equations (E.63.2)--(E.63.3) prove
the desired two-primary vanishing as
well.

We finally identify the adjusted
integral lift. At $(q,\ell)=(3,2)$
the integral local third group is
$\mathbb Z/8$, and at $(2,3)$ it
is $\mathbb Z/3$.
Reduction modulo $N$ is injective
on each, since $8\mid N$ and
$3\mid N$; the kernel of coefficient
reduction is $N$ times the integral
group. Thus the two exceptional
integral obstructions are zero.
B.101 says that this is exactly
the distinguished class $s_*$.

The actual defect of $E_{3,N}\to E_2$
is the reduction of the integral
defect, using the Hall set lift.
Its zero cohomology class provides
continuous finite-valued primitives
for the two central coordinates.
Correcting the set lift by those
primitives gives a continuous
section of this actual quotient.
For an arbitrary prescribed $M$,
choose $K$ with $M\mid K!$ and
project $E_{3,K!}\to E_{3,M}$.
The ordinary pair continues to have
infinite support.
\end{proof}

The modulus in E.63 is fixed before
the pair is constructed. The result
does not give a single pair satisfying
every finite-third-layer condition,
and does not infer continuous integral
third-cohomology vanishing from a
collection of finite reductions.

\begin{result}{E.64}{Proved}{The sharp quadratic-exponent local image needed at the original constants}
\label{coh:sharp-quadratic-image}
Let $n=(j+4)!$, $E(z)=24n(z^2-1)$,
and $D=48n^2$. For $p\mid n$, put
$v=v_p(n)$ and
\[
 F_p(z)=\exp(E(z)\log z).
\]
On the shell
$z\in1+p^v+p^{v+1}\mathbb Z_p$,
its image is exactly
\[
 F_p(1+p^v)(1+p^{W_p}\mathbb Z_p),\qquad
 W_p=
 \begin{cases}
  3v+v_p(48)+1,&p\text{ odd},\\
  3v+7,&p=2.
 \end{cases}
 \tag{E.64.1}
\]
For odd $p$ the map is bijective on that
shell. At two it is two-to-one; restricting
to $z\in1+2^v+2^{v+2}\mathbb Z_2$
makes it bijective onto the same image.
Throughout these shells
$v_p(E(z))=v_p(D)$.

For $p\nmid n$, hence $p\geq5$, on any
residue ball with $z\not\equiv0,\pm1\pmod p$,
the frozen-Teichmuller map
\[
 F_{e,p}(z)=\omega^e
       \exp(24n(z^2-1)\log(z/\omega))
 \tag{E.64.2}
\]
is a bijection onto
$\omega^e(1+p\mathbb Z_p)$, where
$\omega$ is the fixed Teichmuller lift
of that residue. These maps preserve
congruences within their input balls.
For integer inputs whose exponent is
$e$ modulo $p-1$, (E.64.2) is the
ordinary integer power $z^{E(z)}$.
\end{result}
\begin{proof}
Set $z=1+p^vu$. Direct expansion gives
\[
 E(z)\log z=K u^2 H(u),\qquad
 K=48n p^{2v},\qquad
 H(u)=\left(1+\frac{p^vu}{2}\right)
       \frac{\log(1+p^vu)}{p^vu}.
 \tag{E.64.3}
\]
The constant coefficient of $H$ is one,
the linear coefficient is zero, and
its coefficient of $u^k$, for $k\geq1$,
is
\[
       \frac{(-1)^k(k-1)p^{kv}}{2k(k+1)}.
 \tag{E.64.4}
\]
For odd $p$, these nonconstant
coefficients belong to $p\mathbb Z_p$
and tend to zero: for $k\geq2$,
$v_p(k(k+1))\leq k-1$, while the
linear numerator is zero. Consequently
$H\in1+p\mathbb Z_p\langle u\rangle$.
Its square root is an element of
$1+p\mathbb Z_p\langle u\rangle$,
by Hensel lifting in this complete
restricted power-series ring.

At two, $v\geq3$. For $k\geq2$,
$v_2(k)+v_2(k+1)\leq k-1$.
The valuation in (E.64.4) is therefore
at least $k(v-1)\geq4$.
Thus $H=1+16R$ with
$R\in\mathbb Z_2\langle u\rangle$.
The equation
$V+4V^2=R$ has a unique solution in
that ring, by the contraction
$V\mapsto R-4V^2$. Hence
$\sqrt H=1+8V$ is an integral
restricted series.

In both cases the map
$f(u)=u\sqrt{H(u)}$ is a bijection of
$\mathbb Z_p$: it is $u+pG(u)$
for odd $p$ and $u+8G(u)$ at two,
with $G$ integral restricted and
one-Lipschitz. Solving for $u$
is a contraction. The same argument
shows that it maps each residue ball
modulo $p$ onto itself, and at two
each ball modulo eight onto itself.

For odd $p$, on $u\in1+p\mathbb Z_p$
the square $f(u)^2$ ranges bijectively
over $1+p\mathbb Z_p$, since squaring
is a bijection of that group.
Equation (E.64.3) therefore has image
$K(1+p\mathbb Z_p)$ before exponentiation.
The exponential is an isometry on
$p\mathbb Z_p$, and
$v_p(K)=3v+v_p(48)$.
This gives (E.64.1), with its center
equivalently written $F_p(1+p^v)$.

At two, $u$ is odd on the stated
shell. The odd squares are precisely
$1+8\mathbb Z_2$, each with two odd
square roots. Since $f$ preserves
residues modulo eight and is bijective,
the image before exponentiation is
$K(1+8\mathbb Z_2)$, with multiplicity
two. On $u\in1+4\mathbb Z_2$,
$f$ remains a bijection of that
ball and squaring is a bijection
onto $1+8\mathbb Z_2$.
Here $v_2(K)=3v+4$, so the image
after exponentiation is precisely
the coset of depth $3v+7$ in
(E.64.1). The exponential is
used only on $4\mathbb Z_2$.

The valuation of $E(z)$ follows
directly from its two factors.
For odd $p\mid n$, they have
valuations $v$ and zero;
at two they have valuations $v$
and one. Thus in either case
$v_p(24n(z^2-1))=v_p(48n^2)$.

For $p\nmid n$, write
$z=\omega(1+pt)$. Dividing the
logarithm of (E.64.2) by $p$
gives
\[
 \frac{\log(F_{e,p}(z)/\omega^e)}p
   =24n(\omega^2-1)t+pG_0(t),
 \qquad G_0\in\mathbb Z_p\langle t\rangle .
\]
The linear coefficient is a unit.
For $k\geq2$, the logarithmic
coefficient $p^{k-1}/k$ is divisible
by $p$, so the asserted remainder
is integral restricted. The same
contraction gives the stated
bijection, including $p=5,7$.
The frozen exponent $e$ is retained
separately; it is not defined by
evaluating a discontinuous residue
exponent on a variable $p$-adic input.

Finally the logarithm is one-Lipschitz
on every input ball used above,
$E(z)$ is an integral polynomial,
and the exponential is an isometry
on the relevant logarithm ball.
Thus congruence modulo $p^L$ of
the inputs preserves the same
congruence of the outputs.
For an integer input, decomposing
$z=\omega(z/\omega)$ proves the
ordinary-power assertion exactly.
\end{proof}

\begin{result}{E.65}{Proved}{The exact full-class-three valuation criterion and original three-prime recursion}
\label{coh:exact-three-recursion}
Suppose an ordinary completed pair has
valuation patterns $(0,1)$ at two,
$(1,0)$ at three, and at every prime
$q>3$ either $(0,0)$ or one of
\[
 (E_q,0),\quad(0,E_q),\quad(-E_q,-E_q),
 \qquad E_q=24e_q(q^2-1),
 \tag{E.65.1}
\]
where $e_q\geq0$ is an integer and
$e_q\to0$ in $\widehat{\mathbb Z}$.
Then its distinguished integral
class-two class from B.101 lifts
to the actual global integral
class-three quotient. It is the
only class-two class which does.
This proves the valuation-locus
assertion SECTIONS A.58.

Such ordinary infinite-support pairs
exist with exactly the original
parameters of SECTIONS A.59:
\[
 n_j=(j+4)!,\quad D_j=48n_j^2,\quad
 E_j(Q)=24n_j(Q^2-1),
 \tag{E.65.2}
\]
\[
 A=3\prod_{j\geq0}
        \frac{q_j^{E_j(q_j)}}{s_j^{E_j(s_j)}},
 \qquad
 B=-2\prod_{j\geq0}
        \frac{r_j^{E_j(r_j)}}{s_j^{E_j(s_j)}}.
 \tag{E.65.3}
\]
The three sets of primes are pairwise
disjoint and consist of primes greater
than three. The resulting class-three
section is neither a rational-point
image nor a cuspidal image.
\end{result}
\begin{proof}
We first prove the full continuous
obstruction assertion, including
negative support.
Fix $\ell$ and use the actual
class-two representative of B.101.
At a support $q>3$, $q\ne\ell$,
the inertia coordinates $m,n$
are multiples of $E_q$.
The local ordinary cup primitives
in A.25, together with the integral
tame relation, show that the central
inertia coordinate satisfies
\[
                         t\in(E_q/2)\mathbb Z_\ell .
\]
For positive support it is already
a multiple of $E_q$; for the
negative pattern the term
$m(m+1)/2$ accounts for the one
possible denominator two.
There is no further inertia
ambiguity, since an integral
weight-two local cocycle kills
inertia.

Inspect now each term in the two
actual residue polynomials of B.23.
All coefficients other than
the displayed denominators $2,12,24$
are integral at $\ell$:
$q^{-1}$ is a unit,
$(q-1)/2$ is an integer, and
$(q^2-1)/24$ is an integer because
$q>3$ is prime. The binomial
functions at arbitrary integral
coordinates are integer-valued.
The terms $\binom{qm}{2}$ and
$\binom{qn}{2}$ lie in
$(E_q/2)\mathbb Z_\ell$.
Moreover
\[
 B_q(m),B_q(n)\in(E_q/12)\mathbb Z_\ell
\]
by their explicit polynomial formula.
Every remaining monomial contains
one of $m,n,t$.
Consequently
\[
 Z_{D,q},Z_{F,q}\in(E_q/24)\mathbb Z_\ell,
 \qquad
 \frac{Z_{D,q}}{q^2-1},
 \frac{Z_{F,q}}{q^2-1}\in e_q\mathbb Z_\ell .
 \tag{E.65.4}
\]
This checks the full formulas and
does not require cancellation
specific to positive support.
At a zero-support prime, $m=n=t=0$,
so both residues are zero.
At the two old primes, the only
possibly nonunit residue moduli
are $(q,\ell)=(3,2),(2,3)$;
B.101 makes both obstruction
classes zero. Their finitely
many normalized residues are
therefore integral.
Equation (E.65.4) and $e_q\to0$
make the entire normalized
sequences integral and $c_0$.
A.22 and B.23, including the
zero fixed-ramification kernel
in weight three at two, prove
that the actual continuous
global obstruction is zero.
Continuous primitives give
the integral class-three section.
Uniqueness of its class-two
projection is B.101.

It remains to establish the exact
arithmetic construction, since
E.31's general bound $J\geq2d^2$
would not by itself permit $J=4$
when $d=2$. Keep the original
parameters (E.65.2) and define
\[
 B(D)=(D+1)^4,\quad
 d_j(p)=
 \begin{cases}
  2v_p(n_j)+v_p(96),&p\mid n_j,\\
  0,&p\nmid n_j,
 \end{cases}
 \quad H_j(p)=2d_j(p)+j+3 .
 \tag{E.65.5}
\]
Maintain pairwise coprime coefficients
$a=3u^{24},b=-2v^{24},c=w^{24}$
with disjoint positive supports,
unit-coordinate solubility of
\[
 aX^{D_j}+bY^{D_j}=cZ^{D_j}
                         \quad\text{over every }\mathbb Q_p,
 \tag{E.65.6}
\]
and
$v_p(a+b-c)\geq H_j(p)$ for
$p\leq B(D_j)$.
The initial triple $(3,-2,1)$
satisfies them exactly.

The three radical fields used at
stage $j$ are precisely those of
(E.31.6) with $n=n_j$.
Their maximal abelian subfield
is $\mathbb Q(\mu_n)$.
That proof used only $24\mid n$
and the coefficient forms just
specified, not the larger
general initial factorial.
It includes the third square
class six, the denominator in
$-a/b$, and the nonabelian
pure-quartic subfield.
Thus each prescribed unit CRT
class equal to one modulo $n$
can be imposed simultaneously
with splitting in its respective
radical field.

Choose $P\geq B(D_{j+1})$ at
least every old support prime.
Before processing primes
$p\leq P$ in increasing order,
choose precisions
$L_p\geq H_{j+1}(p)$,
$p^{L_p}>P$, and
$L_p\geq v_p(n)+2$ when $p\mid n$.
For each of the three future
prime variables impose
\[
 \begin{cases}
  v_p(z-1)=v_p(n),&p\mid n,\\
  z\not\equiv0,1,-1\pmod p,&p\nmid n.
 \end{cases}
 \tag{E.65.7}
\]
The second case occurs only at
$p\geq5$, where permitted residues
exist, including $p=5,7$.

When $p$ is reached, all prime
factors of $p-1$ have already
been processed with enough
precision to fix each input
modulo $p-1$. They therefore
fix separately the three exponent
classes $e_i=E_j(z_i)\bmod(p-1)$.
For every prime $\lambda\mid p-1$,
(E.65.7) gives the exact valuation
\[
 v_\lambda(E_j(z_i))
                 =v_\lambda(48n^2).
 \tag{E.65.8}
\]
For odd $\lambda\mid n$ this uses
the valuations $v_\lambda(n),0$
of $z_i-1,z_i+1$; at two it uses
$v_2(n),1$. If $\lambda\nmid n$,
both sides are zero by exclusion
of the two residues.
Consequently the three finite-field
power maps have exactly the same
image as the $D_j$th-power map.
Both forbidden residues $1,-1$
map to one, since the exponents
are even.

At $p\mid n$, use E.64. Its
precise image depths satisfy
$W_p\leq H_j(p)$. In fact the
respective values of $H_j$ are
$4v+j+13$ at two, $4v+j+5$
at three, and $4v+j+3$ at the
other coefficient primes; the
values of $W_p$ are $3v+7$,
$3v+2$, and $3v+1$.
At two use the smaller shell
from E.64 if a single-valued
inverse is desired.
All three output maps have the
same full coset. At their common
base value the linear combination
has error $(a+b-c)F_0$, divisible
by $p^{W_p}$. Vary an output
whose coefficient is a unit to
cancel that error exactly.

At $p\nmid n$ use the three
frozen-exponent bijections
(E.64.2). If $p\leq B(D_j)$,
choose a common finite-field
$D_j$th-power value $\lambda$.
When the image group is nontrivial
take $\lambda\ne1$, so every
preimage avoids $1,-1$. When
the image is trivial all nonzero
residues are preimages, and
$p\geq5$ still leaves permitted
ones. The three output balls
are the same Teichmuller coset.
The invariant $a+b-c\equiv0\pmod p$
and a unit coefficient permit
an exact correction inside it.
If instead $p>B(D_j)$, start
with (E.65.6) and scale all three
unit coordinates by one unit.
The group of $D_j$th powers has
more than three elements, whereas
at most three scaling values
make one of the resulting power
values equal to one.
Choose any other scaling.
The three exact unit outputs
then have permitted residue
preimages, and E.64 lifts them
to the required local inputs.
This covers large old support
primes as well.

Retain the constructed inputs
modulo $p^{L_p}$.
Congruence preservation in E.64
retains error at least
$H_{j+1}(p)$; earlier residues
also continue to fix the
later frozen exponents.
CRT gives three unit classes
modulo $\prod_{p\leq P}p^{L_p}$,
each one modulo $n$.
Use the radical compatibility
to choose distinct new primes
$q,r,s>P$ in these classes,
splitting the respective old
radical fields. Set
\[
 a'=a q^{E_j(q)},\quad
 b'=b r^{E_j(r)},\quad
 c'=c s^{E_j(s)}.
 \tag{E.65.9}
\]
Every new exponent is divisible
by $24$, so the coefficient
forms and disjoint supports
are preserved.

At $p\leq P$, the new error is
at least $H_{j+1}(p)$ and
$H_{j+1}(p)>2v_p(D_{j+1})$.
The one-variable strong Hensel
estimate, applied to a variable
with unit coefficient at
$(1,1,1)$, gives the next
unit-coordinate point.
At a new support prime
$Q\in\{q,r,s\}$, (E.65.7) gives
the stronger exact identity
\[
                   \gcd(D_{j+1},Q-1)=n_j.
 \tag{E.65.10}
\]
Indeed primes dividing $n_j$
have exactly that valuation in
$Q-1$; any new prime of
$D_{j+1}$ is at most $j+5\leq P$
and was required not to divide
$Q-1$. Also $n_j\mid D_{j+1}$.
Thus the $n_j$th-power and
$D_{j+1}$th-power images coincide.
At $q$ the split ratio $c/b$
and the new $r,s$ factors make
$c'/b'$ such a power.
With $X=Z=1$, simple-root
Hensel solves
$Y^{D_{j+1}}=(c'-a')/b'$.
At $r$, use the ratio $c/a$
to solve for $X$ with $Y=Z=1$.
At $s$, use $-b/a$, the inverse
of the third split radicand,
to solve
$X^{D_{j+1}}=(c'-b')/a'$.
No congruence between two newly
chosen primes is needed.

At every remaining prime all
coefficients are units and
$p>B(D_{j+1})$. The nonzero
finite-field solution bound
proved in E.102, or the explicit
Weil-bound calculation in E.31,
gives a unit point and simple
Hensel lifts it. This completes
both all-place invariants.

The limit argument is literal.
Every new exponent is divisible
by the growing factorial $n_j$,
so all three coefficient products
converge in the global power
completion. At a fixed finite
prime their valuations stabilize
and their unit factors converge,
giving nonzero ordinary limits
$a_p,b_p,c_p$.
The increasing precision
$H_j(p)$ eventually applies
and tends to infinity, so
$a_p+b_p=c_p$. Their quotients
$x_p=a_p/c_p$,
$1-x_p=b_p/c_p$ are the required
ordinary local point coordinates.
The valuation patterns are
exactly (E.65.1), with
$e_{q_j}=e_{r_j}=e_{s_j}=n_j$.
For any fixed modulus only
finitely many of these fail
to be divisible by it; thus
$e_q\to0$ in $\widehat{\mathbb Z}$.
The first part of the proof
supplies the full integral
class-three section.

Infinite support excludes any
rational-point pair by E.101.
The cusp-zero abelian image has
trivial second coordinate,
contradicting $v_2(B)=1$;
cusp one similarly contradicts
$v_3(A)=1$; and the two
valuations of the cusp-infinity
image are equal, contradicting
the pattern already at two.
These are statements about the
actual abelian cusp maps and
remain valid in class three.
\end{proof}

The supplied A.59 text included an
unavailable coefficient-prime lemma.
E.64 proves every local-image
property required by its recursion,
including the exact two-adic image
and the small coefficient primes.
No missing lemma is being treated
as an established input.
